\section{本章小结}
\label{sec:mmap-summary}

本章开头的问题是：输入的误差经过映射后变成多大的输出误差，近似映射的极限保留什么性质，反复代入何时能求出方程的解。第一层回答是连续性本身。在度量空间之间，第~\ref{chap:topological-spaces}章的开集原像定义等价于逐点的$\varepsilon$-$\delta$保证，也等价于保持收敛点列的极限；连续性保证原像保持开闭，不保证像保持开闭；序列判据的反向依赖每点有可数个球作邻域基，一般拓扑空间中只有正方向成立。

第二层回答是统一的误差控制。均匀连续要求阈值与位置无关，从而把柯西列送到柯西列；Lipschitz连续进一步要求输出误差不超过输入误差的固定倍数，多步计算的放大倍数至多相乘。三种条件依次加强，$x^2$与$\sqrt x$分别说明连续不必均匀连续、均匀连续不必Lipschitz连续。紧定义域上连续自动均匀连续，等度连续族自动一致等度连续，但紧性不提供Lipschitz常数。

第三层回答涉及映射列与映射族。逐点收敛允许达到精度所需的序号随输入改变，一致收敛则要求统一的序号；一致极限保持连续，而这只需目标空间是度量空间。目标空间完备时，一致柯西列必有一致极限，由此有界映射空间与有界连续映射空间在一致度量下完备，这是后面函数空间完备性的共同依据。等度连续要求全族共用阈值，为第~\ref{chap:function-class-covers}章的紧性判据作准备。

最后，映射条件与空间的性质配合使用。等距映射保留一切度量性质；同胚只保留拓扑性质，实轴与开区间说明完备、有界都不是拓扑性质；同一集合上的度量按拓扑等价、一致等价、强等价分别保留收敛与紧性、柯西列与完备性、有界性与误差尺度。完备目标空间中，稠密子集上的均匀连续映射唯一地延拓到全空间，并保留阈值与Lipschitz常数，由此完备化在等距意义下唯一。压缩映射定理把统一的压缩常数与完备性结合起来：几何级数给出柯西列，完备性给出极限，连续性证明它是不动点，先验与后验误差界则告诉我们何时可以停止迭代。
